% from ega0-18-fr.tex
\oldpage[0\textsubscript{IV-1}]{68}

continu de $B$-bimodules $L\to L'$, on en déduit canoniquement un
homomorphisme de groupes additifs
\[
  \operatorname{Exantop}_A(B,L)\longrightarrow
  \operatorname{Exantop}_A(B,L')
\]
par passage à la limite inductive à partir de (18.3.6.1).
\end{env}

\begin{env}[18.5.3]
\label{0.18.5.3-fr}
Étant donnés un anneau topologique $C$ et deux $C$-bimodules topologiques
$M$, $N$,\linebreak on note $\operatorname{Hom.\ cont}_C(M,N)$ le groupe additif des
$C$-homomorphismes \emph{continus} de $M$ dans $N$.
\end{env}

\begin{lemma}[18.5.3.1]
\label{0.18.5.3.1-fr}
Soient $C$ un anneau topologique, $E$, $L$ deux $C$-bimodules topologiques~;
on suppose que les topologies sont linéaires, que $L$ est discret et annulé
par un idéal bilatère ouvert de $C$. Alors on a un isomorphisme canonique
\begin{equation}
\label{0.18.5.3.2-fr}
\tag{18.5.3.2}
  \varinjlim\operatorname{Hom}_{C/\mathfrak R}(E/V,L)
  \xrightarrow{\sim}\operatorname{Hom.\ cont}_C(E,L)
\end{equation}
où dans le premier membre la limite inductive est prise suivant l'ensemble
ordonné filtrant des couples $(\mathfrak R,V)$ tels que $\mathfrak R$ soit un
idéal bilatère ouvert de $C$, $V$ un sous-$C$-bimodule ouvert de $E$, tels que
$\mathfrak R.L=L.\mathfrak R=0$, $\mathfrak R.E\subset V$,
$E.\mathfrak R\subset V$.
\end{lemma}

\begin{proof}
Comme $C/\mathfrak R$ et $E/V$ sont discrets, on a des homomorphismes
canoniques
\[
  w_{\mathfrak R,V}:\operatorname{Hom}_{C/\mathfrak R}(E/V,L)
  \longrightarrow\operatorname{Hom.\ cont}_C(E,L)
\]
formant un système inductif, d'où un homomorphisme (18.5.3.2) par passage à la
limite inductive. Comme l'homomorphisme $E/V'\to E/V$ est surjectif pour
$V'\subset V$, il résulte aussitôt de la définition que l'homomorphisme
\[
  \operatorname{Hom}_{C/\mathfrak R}(E/V,L)
  \longrightarrow\operatorname{Hom}_{C/\mathfrak R}(E/V',L)
\]
(avec $\mathfrak R.L=L.\mathfrak R=0$, $\mathfrak R.E\subset V'$,
$E.\mathfrak R\subset V'$) est injectif, et il en est évidemment de même de
l'homomorphisme
\[
  \operatorname{Hom}_{C/\mathfrak R}(E/V,L)
  \longrightarrow\operatorname{Hom}_{C/\mathfrak R'}(E/V,L)
\]
pour $\mathfrak R'\subset\mathfrak R$~; on en conclut que l'homomorphisme
(18.5.3.2) est injectif. D'autre part, si $u$ est un $C$-homomorphisme continu
de $E$ dans $L$, son noyau est un sous-bimodule ouvert $V_0$ de $E$, et si
$\mathfrak R_0$ est un idéal bilatère ouvert de $C$ tel que
$\mathfrak R_0.L=L.\mathfrak R_0=0$ et $\mathfrak R_0.E\subset V_0$,
$E.\mathfrak R_0\subset V_0$, il est clair que $u$ est l'image canonique d'un
$(C/\mathfrak R_0)$-homomorphisme de $E/V_0$ dans $L$, donc (18.5.3.2) est
surjectif.
\end{proof}

Cela étant, la proposition (18.3.8) se généralise comme suit aux anneaux
topologiques~:

\begin{proposition}[18.5.4]
\label{0.18.5.4-fr}
Soient $B$ un anneau topologique linéairement topologisé, $\mathfrak J$ un
idéal bilatère de $B$, $C=B/\mathfrak J$ l'anneau topologique quotient~;
$\mathfrak J/\mathfrak J^2$ (où $\mathfrak J$ est muni de la topologie induite
par celle de $B$ et $\mathfrak J/\mathfrak J^2$ de la topologie quotient de
celle de $\mathfrak J$) est alors canoniquement muni d'une structure de
$C$-bimodule topologique. Pour tout $C$-bimodule discret $L$ annulé par un
idéal ouvert de $C$, il existe alors un isomorphisme canonique
\begin{equation}
\label{0.18.5.4.1-fr}
\tag{18.5.4.1}
  \operatorname{Hom.\ cont}_C(\mathfrak J/\mathfrak J^2,L)
  \xrightarrow{\sim}\operatorname{Exantop}_B(C,L).
\end{equation}
\end{proposition}

\begin{proof}
En effet, pour tout idéal ouvert $\mathfrak R$ de $B$ tel que
$(\mathfrak J+\mathfrak R)/\mathfrak J$ annule $L$, on a, d'après (18.3.9),
un isomorphisme canonique
\[
  \operatorname{Hom}_{B/(\mathfrak J+\mathfrak R)}
    ((\mathfrak J+\mathfrak R)/(\mathfrak J^2+\mathfrak R),L)
  \xrightarrow{\sim}
  \operatorname{Exan}_{B/\mathfrak R}(B/(\mathfrak J+\mathfrak R),L)
\]
et il suffit de passer à la limite inductive, en tenant compte de (18.5.3.1).
\end{proof}

% (continued in ega0-19-fr.tex, before its first page marker)
\section{Algèbres formellement lisses et anneaux de Cohen}
\label{section:0.19-fr}
\label{0.19-fr}

\setcounter{subsection}{-1}
\subsection{Introduction}
\label{subsection:0.19.0-fr}
\label{0.19.0-fr}

